% GENERATED by convert_sources.py from docs/requirements.md. Do not edit by hand.
\subsection{Portable Stereo Bluetooth Speaker --- Requirements (v0.1)}

Working name: \textbf{TBD}. Goal: a ``vibe hardware'' mixed-signal project (digital + analog + power) that holds its own against a JBL Flip/Charge-class speaker, built from a JLCPCB-assembled PCB, LCSC parts, 3D-printed enclosure (Flashforge Adventurer 5M Pro), and off-the-shelf drivers/hardware.

\subsubsection{1. Product targets}

{\def\LTcaptype{none} % do not increment counter
\begin{longtable}[]{@{}
  >{\raggedright\arraybackslash}p{(\linewidth - 4\tabcolsep) * \real{0.3333}}
  >{\raggedright\arraybackslash}p{(\linewidth - 4\tabcolsep) * \real{0.3333}}
  >{\raggedright\arraybackslash}p{(\linewidth - 4\tabcolsep) * \real{0.3333}}@{}}
\toprule\noalign{}
\begin{minipage}[b]{\linewidth}\raggedright
Item
\end{minipage} & \begin{minipage}[b]{\linewidth}\raggedright
Target
\end{minipage} & \begin{minipage}[b]{\linewidth}\raggedright
Notes
\end{minipage} \\
\midrule\noalign{}
\endhead
\bottomrule\noalign{}
\endlastfoot
Channels & Stereo, 2-way per side, \textbf{bi-amped} (active DSP crossover) & 4 driver channels total \\
Drivers & 2× woofer (2.5--3.5''), 2× tweeter (\textasciitilde¾''), 2× passive radiator & PRs on end caps (visible) \\
Bass extension & F3 ≤ 65 Hz (with DSP), F6 ≤ 55 Hz & JBL Flip 6 ≈ 63 Hz \\
Max SPL & ≥ 90 dB @ 1 m (pink noise, battery, \textbf{on a table / half-space}) & limiter-protected; free-air ≈ 85 dB \\
Driver layout & Woofers rear-mounted on front face (66 mm apart); \textbf{tweeters angled up 40° in sealed printed pods} on the front shoulders (166 mm apart), \textbf{minimal 3 mm bezel instead of a full waveguide (decided; verify with BEM sim)} & see \texttt{docs/\allowbreak{}mech\_\allowbreak{}envelope.md} \\
Enclosure (decided) & \textbf{Ø95 × 210 mm body, 216 mm over PR flanges}, 5 printed PETG parts, \textasciitilde0.85 L net & nothing in stock fits Ø90 \\
Battery pick (decided) & \textbf{Gens ace 2S 5000 mAh TX/RX LiPo} (109×42×19.5 mm, \textasciitilde\$31), bonded 10 k NTC; fallback Gens ace 3500 mAh with spacer & RC-brand, no PCM: board protection covers it; bench-test 6 A \\
Amp power & \textasciitilde15--20 W/woofer, \textasciitilde5--10 W/tweeter & boost rail \textasciitilde18--20 V \\
Source & Bluetooth A2DP sink (SBC required, AAC nice-to-have) & \\
Control & Wi-Fi web UI for live EQ/presets; rotary encoder + push; 2--3 buttons & \\
Display & 1.28'' round TFT (GC9A01-class) on top & mode, EQ curve, level meter, battery \\
Mic & 1× digital MEMS mic & placement (bass) compensation + ambient auto-loudness \\
Battery & Off-the-shelf \textbf{pre-built 2S pack} (no spot welding), plugs in via \textbf{XT30 (power) + JST-XH 3-pin (balance)}; charger/balancing/protection stay on our board. Pouch pack preferred; \textbf{no device-grade pack fits the capsule, so the decided pack is the Gens ace 2S 5000 mAh TX/RX LiPo} (see Battery pick row); enclosure needs swell clearance + puncture protection; charge current ≤ 0.5C unless the pack datasheet says otherwise & JLC won't ship cells; ≤ \textasciitilde40 W avg draw (charger BATFET 6 A RMS) \\
Drivers (\textbf{set B8, decided}) & 2× Dayton ND65-8 woofer (8 Ω), 2× Peerless OT19NC00-04 tweeter (4 Ω), 2× Dayton DMA58-PR & \$93.54 at Parts Express; F3 64 Hz, \textasciitilde90.6 dB pink (half-space, A1 refined model), ≥0.50 L net; see \texttt{research/\allowbreak{}drivers.md} \\
Amp rail & \textbf{PVDD 18 V} (TPS61178), DNP trim for dynamic PVDD & 8 Ω woofers keep the TAS5805M well under its 5 A OCP \\
Runtime & ≥ 8 h at moderate volume & \\
USB & USB-C: PD sink charging/power, play-while-charging, firmware flashing & USB-audio input = \textbf{optional add-on, decide after research} \\
Size & Horizontal capsule, ≤ \textasciitilde210 mm long, \textasciitilde80--95 mm dia & printer volume 220×220×220 mm \\
Budget & \textasciitilde\$250 for v1 (5 PCBA, drivers, filament, hardware, UMIK-1 mic) & \\
\end{longtable}
}

Out of scope: waterproofing, multi-speaker linking, phone app (web UI instead), LDAC/aptX.

\subsubsection{2. Block diagram}

{\scriptsize
\begin{verbatim}
 USB-C ──► PD sink ──► 2S charger + power path ──► VSYS (6–8.6 V)
              │               │                       │
              │           2S Li-ion pack ◄── protection/balancing
              │                                       ├──► boost ~18–20 V ──► AMP_L (TAS5805M-class) ──► woofer L / tweeter L
              │                                       │                  └──► AMP_R                ──► woofer R / tweeter R
              │                                       └──► buck 3.3 V ──► ESP32 + UI + mic
              └──► USB-UART (flashing/debug) ──► ESP32-WROVER-E (PSRAM required for BT+Wi-Fi audio)
                                                  ├─ I2S1 TX ──► both amps (I2S)
                                                  ├─ I2S0 PDM RX ◄── MEMS mic (PDM, on UI board)
                                                  ├─ I2C ──► amps (DSP coeffs), charger, fuel gauge
                                                  ├─ SPI ──► round TFT
                                                  └─ GPIO ◄── encoder, buttons, battery sense
\end{verbatim}
}

\subsubsection{3. DSP split}

\begin{itemize}
\tightlist
\item
  \textbf{Amp DSP (TAS5805M-class):} crossover, per-driver EQ, DRC/limiter, excursion protection. Safe even if MCU firmware glitches.
\item
  \textbf{ESP32:} A2DP decode, user EQ/presets, loudness, bass enhancement, stereo widening, mic-based auto-EQ, web UI, display.
\end{itemize}

\subsubsection{4. Manufacturing / mechanical}

\begin{itemize}
\tightlist
\item
  PCB: 4-layer, JLCPCB PCBA, LCSC parts, prefer JLC \textbf{basic} parts; minimize extended-part count. Main board + possibly small top UI board.
\item
  Enclosure: FDM PETG/ASA, ≥4 perimeters, dense infill, internal bracing, printed tweeter waveguides, foam gaskets, brass heat-set inserts, perforated-aluminum grille, polyfill stuffing.
\item
  Hardware from McMaster-Carr / Home Depot.
\end{itemize}

\subsubsection{5. Simulation \& verification plan}

{\def\LTcaptype{none} % do not increment counter
\begin{longtable}[]{@{}
  >{\raggedright\arraybackslash}p{(\linewidth - 4\tabcolsep) * \real{0.3333}}
  >{\raggedright\arraybackslash}p{(\linewidth - 4\tabcolsep) * \real{0.3333}}
  >{\raggedright\arraybackslash}p{(\linewidth - 4\tabcolsep) * \real{0.3333}}@{}}
\toprule\noalign{}
\begin{minipage}[b]{\linewidth}\raggedright
Domain
\end{minipage} & \begin{minipage}[b]{\linewidth}\raggedright
Our own tool
\end{minipage} & \begin{minipage}[b]{\linewidth}\raggedright
Cross-check
\end{minipage} \\
\midrule\noalign{}
\endhead
\bottomrule\noalign{}
\endlastfoot
Box/PR alignment & Python lumped-element (T/S) model & WinISD/VituixCAD (Wine) if needed \\
Directivity/diffraction & --- & bempp-cl + gmsh on FreeCAD geometry \\
PCB thermal & Python 2D finite-difference solver + heatmaps & FreeCAD FEM (CalculiX) / Elmer \\
Acoustic tuning & --- & REW (macOS) + UMIK-1, measured → amp EQ \\
\end{longtable}
}

\subsubsection{6. Environment}

macOS arm64, KiCad + FreeCAD installed, Python 3.9 (system), Flashforge Adventurer 5M Pro on network.

\subsubsection{7. Open decisions}

\begin{itemize}
\tightlist
\item
  USB-audio input (extra bridge chip) --- yes/no after research.
\item
  Boost rail vs 3S pack vs direct-from-battery amp supply.
\item
  Driver set (sets enclosure volume and size).
\end{itemize}
